\documentclass{article}

% Language setting
% Replace `english' with e.g. `spanish' to change the document language
\usepackage[english]{babel}

% Set page size and margins
% Replace `letterpaper' with `a4paper' for UK/EU standard size
\usepackage[a4paper,top=2cm,bottom=2cm,left=3cm,right=3cm,marginparwidth=1.75cm]{geometry}

% Useful packages
\usepackage{amsmath}
\usepackage{graphicx}
\usepackage[colorlinks=true, allcolors=blue]{hyperref}

\title{Discrete Math Course}
\author{Wisam Yassar Mahardika}
\date{April 4\textsuperscript{th} 2026}

\begin{document}
\maketitle
\clearpage
\tableofcontents
\clearpage
\section{Propositions, Negations, Conjunctions and Disjunctions}
\subsection{Propositions}
A \textbf{Proposition} is a declarative statement that is either true or false.
\begin{itemize}
\item The sky is blue
\item The moon is made of cheese
\item Luke, I am your father
\item Sit down
\item $ X + 1 = 2 $ \\
\end{itemize}

 \textit{Sit down is neither true or false so it is not a proposition text} \\
\textit{$ X + 1 = 2 $ is also neither true or false because we don't know the value of x}
\\

\subsection{Constructing Compound Propositions}
A \textbf{Compound Proposition} is comprised of propositions and one or more of the following connectives:
\\


\begin{itemize}
\item Negation $ \neg $ ``not`
\item Conjunction $ \land $     ``and``
\item Disjunction $ \lor $ ``or``
\item Implication $ \implies $ ``if, then``
\item Biconditional $ \Longleftrightarrow{} $ ``if and only if``\\
\end{itemize}

Each proposition is represented by a propositional variable \textit{p, q, r, s ...}.\\


Example: $ p \implies q $

\subsubsection{Negation}
The \textbf{negation} of the proposition $ p $ is $ \neg p $ (not $p$) \\


Example if $ p $ denotes ``The grass is green``, then $ \neg p $ denotes ``it is not the case that the grass is green`` or more simply ``the grass is not green``.\\


\begin{itemize}
\item My dog is the cutest dog  ($ \neg p $ My dog is \textbf{not} the cutest dog)
\item The door is not open ($ \neg p $ The door \textbf{is} open)
\item Are we there yet
\end{itemize}
\textit{The third statement is not a proposition because it is nor true or false.}

\clearpage

\subsubsection{Truth Table}
Each row of a \textbf{truth table} gives us one possibility for the truth values of our proposition(s), since each proposition has two possible truth value, \textit{true} or \textit{false}, we will have 2 rows for each propositions (or $ 2^n $ rows where $ n = \# $ of propositions). \\


\textit{Truth Table for $ \neg p $} 
\begin{tabular}{|c|c|}
\hline
$ p $ & $ \neg p $ \\
\hline
T & F \\
F & T \\
\hline
\end{tabular} \\
\textit{$ p $ is combinations and $ \neg p $ is connectives} \\


\subsubsection{Conjunction}
The \textbf{conjunction} of propositions $ p $ and $ q $ is denoted  $ p \land q $ and read $ p $ ``and`` $ q $.
For a conjunction to be true, \textbf{both} propositions must be true. \\

\textit{Truth Table for $ p \land q $} 
\begin{tabular}{|c|c|c|}
\hline
$ p $ & $ q $ & $ p \land q $ \\
\hline
T & T & T \\
T & F & F \\
F & T & F \\
F & F & F \\
\hline
\end{tabular}\\


\subsubsection{Disjunction}
The \textbf{disjunction} of propositions $ p $ and $ q $ is denoted $ p \lor q $ and read $ p $ ``or`` $ q $.
For a disjunction to be true. \textbf{either} proposition must be true.\\


\textit{Truth Table for $ p \lor q $}
\begin{tabular}{|c|c|c|}
\hline
$ p $ & $ q $ & $ p \lor q $ \\
\hline
T & T & T \\ 
T & F & T \\
F & T & T \\
F & F & F \\
\hline
\end{tabular}\\


\subsubsection{The Connective ``or`` in English ``xor``}
\begin{itemize}
	\item \textbf{Inclusive ``or`` $ p \lor q $}
	The prerequisite for MA420 is either MA315 or MA335.
	\item \textbf{Exclusive ``or`` $ p \oplus q $}
	You get soup or salad with your entree.
\end{itemize}


\textit{Truth Table for $ p \oplus q $}
\begin{tabular}{|c|c|c|}
\hline
$ p $ & $ q $ & $ p \oplus q $ \\
\hline
T & T & F \\ 
T & F & T \\
F & T & T \\
F & F & F \\
\hline
\end{tabular}\\


\end{document}
